\begin{proposition}[Period-invariant fans after base change]\hypertarget{tag-0302}{}\label{push:src:B.2.1}\label{push:B21}
\StatementRecord{push:src:B.2.1}{proved}{0302}{\SeriesRef{push:src:A.2.1}, \SeriesRef{push:src:A.1.2}, \SeriesRef{push:src:A.2.2}}

For the plain base change $u=r^n$ of the untwisted cusp
model, the smooth period-invariant refinements which
retain the prescribed bottom triangulation, are jointly
flat over $(q,r)$, and have reduced $r$-central divisor
are indexed by $k=0,\ldots,n$. In a fundamental prism
write its four bottom rays as $A,B,C,D$ as in \ref{push:src:A.2.1},
and put $\rho_i=(0,i,0,1)$ for $0\leq i\leq n$.
The maximal cones for the index $k$ are
\[
 ABC\rho_k,\quad BCD\rho_k,\quad
 AB\rho_i\rho_{i+1}\ (i<k),\quad
 CD\rho_i\rho_{i+1}\ (i\geq k).
\]
There are thirteen such models for $n=12$ and three
for $n=2$. Each has smooth total space and a reduced
normal-crossings central divisor with $n$ components
and no triple locus. Its component labels over the
punctured $q$-disc are globally $\mathbb Z/n$; in
particular their monodromy around $q=0$ is the identity.
\end{proposition}


\begin{proof}
\begin{enumerate}[label=\arabic*.,leftmargin=*]
\item The map of cocharacter lattices for the base change is
\[
 (y_1,y_2,h_q,h_r)\longmapsto
 (y_1,y_2,h_q,nh_r).
\]
The old top ray indexed by $l$ pulls back to
$(0,nl,0,1)$. Smoothness on the two-dimensional top
face forces insertion of every primitive ray
$\rho_{nl+i}$, $0<i<n$: a skipped interval of length
$d>1$ would retain lattice index $d$ on that face.
Joint flatness excludes any new ray with both base
heights positive, and the fixed bottom excludes new
rays of $h_r=0$. Reducedness permits only top rays
of $h_r=1$. These are consequently all the rays.

\item On the left face $AC\rho_0\rho_n$, a triangulation
with these collinear top vertices has a unique triangle
containing the bottom edge $AC$. Write its third
vertex as $\rho_a$. The triangles on either side
of that bridge are then forced. The analogous right
face has bridge $BD\rho_b$. The first period takes
$A,C$ to $B,D$ and fixes every top ray, so periodicity
requires $a=b=k$. The fixed bottom triangles $ABC$
and $BCD$ must then cone to $\rho_k$. The remaining
successive top intervals force exactly the displayed
cones, leaving no choice. Conversely those cones
have determinant $\pm1$ and common-face
intersections. Their normalized volumes sum to
$n+2$, the volume of the prism, so they cover it.
For $n=2$ and $n=12$, the determinant, intersection and
volume calculations use exact rational arithmetic.

\item The period actions are
\[
 G_1:(y_1,y_2)\mapsto(y_1+h_q,y_2),\qquad
 G_2:(y_1,y_2)\mapsto(y_1,y_2+h_q+nh_r).
\]
Thus the choices propagate consistently from the
fundamental prism. The proper toric subdivisions
descend by the covering construction of \ref{push:src:A.1.2}.
Every ray has pure $q$-height or pure $r$-height, so
in a regular affine chart $q$ and $r$ are monomials
in disjoint sets of coordinates. This proves joint
flatness. All nonzero $r$-heights equal one, and
there are at most two top rays in any maximal cone;
this proves reduced normal crossings and absence of
triple points. Properness follows by descent from
the proper toric modification of the base-changed
proper cusp model.

\item The central components and all their intersection
maps are specified by the stars of the $\rho_i$.
Modulo the periods their indices are $i\bmod n$;
adjacent indices meet in the conductor elliptic
surface. Away from the cusp their normalizations
are the ruled bundles of the expanded plumbing.
At the unique bridge class the star has both bottom
triangles, and gives the nodal-curve blow-up and
small resolution described in \ref{push:src:A.2.2}, on the
corresponding ruled bundle. This star description
also fixes the gluing when the bridge is on an
inserted component. The endpoint choices $k=0,n$
are the canonical resolutions of the base changes
of the two original models. Intermediate values
are related to them by the successive side flops.

\item Finally, $G_1$ fixes each index and $G_2$ changes it
by $n$. The units in the cusp periods are
single-valued holomorphic functions of $q$ and do
not change a ray label. The toric charts themselves
are defined over $q$, not a root of $q$. Hence
these labels remain global on the punctured
$q$-disc, proving the monodromy assertion. Further
smooth periodic subdivisions with mixed rays give
infinitely many models flat over $r$ alone, as in
\ref{push:src:A.2.1}; they are excluded by the stated joint-flatness
condition.
\end{enumerate}
\end{proof}
